\documentclass[10pt,a4paper]{amsart}
\usepackage[margin=19mm]{geometry}
\usepackage{amsmath,amssymb,mathtools}
\usepackage[scr=rsfs,cal=euler]{mathalfa}
\usepackage{xfrac}
\usepackage[unicode,hidelinks,bookmarks=false]{hyperref}
\newcommand{\ie}{i.e.\ }
\newcommand{\cf}{cf.\ }
\newcommand{\resp}{respectively\ }
\newcommand{\Subsection}{\subsection}
\providecommand{\nobreakdash}{\nobreak\hbox{-}}
\setlength{\emergencystretch}{2em}
\multlinegap0pt
\usepackage{amssymb}
\usepackage{tikz}
\usepackage{csquotes}
\usepackage{mathtools}
\newtheorem{lemma}{Lemma}
\newtheorem{theorem}{Theorem}
\newcommand{\Dyck}{\mathbb{D}}
\DeclareMathOperator{\area}{\mathrm{area}}
\DeclareMathOperator{\touch}{touch}
\title{Shalika germs for tamely ramified elements in $GL_n$}
\author{Oscar Kivinen, Cheng-Chiang Tsai}
\date{Source: arXiv:2209.02509v4 (CC BY 4.0)}
\begin{document}
\maketitle

\noindent\emph{Source and licence.} Shalika germs for tamely ramified elements in $GL_n$. Version arXiv:2209.02509v4 (CC BY 4.0), distributed by arXiv under the Creative Commons Attribution 4.0 International licence, http://creativecommons.org/licenses/by/4.0/, as recorded in the arXiv licence metadata for this exact version and shown on the abstract page https://arxiv.org/abs/2209.02509v4. Full licence notice: \texttt{licenses/arxiv-2209.02509-CC-BY-4.0.txt}.\par\smallskip

\noindent\textit{Editorial scope.} Counted unit: the article's theorem thm:comparerecursion (source0.tex lines 1976--1982). The article's complete original proof of it is delivered with it (source0.tex lines 1983--2020, one \texttt{proof} environment of 38 lines). No mathematics is written, repaired or re-derived: the statement and the proof are the article's own lines, in the article's own notation, and no retained line is substituted or re-flowed.  Retained uncounted as the source-local context the proof needs: lines 1878--1881; lines 1937--1943; lines 1969--1973.  Nothing was omitted from any retained span, so the package carries no omission record.  No retained line contains a citation, so the wrapper carries no bibliography and no bibliography entry.  No retained line contains a figure environment, an image-inclusion command or drawing code, so no image is delivered and the package declares no asset dependency.  Editorial formatting only, in addition: an AMS \texttt{amsart} preamble that restates the article's own declarations and macros used by the retained lines, namely the environments \texttt{lemma}, \texttt{theorem} on separate counters, together with the standard packages those lines need.  The retained environments print in this document as Lemma 1, Theorem 1 in order of appearance, whereas the article prints the same items with the numbering of its own full document.  The complete native source of the article as distributed in the arXiv e-print for this exact version is bundled unchanged as \texttt{sources/130349/source0.tex}.
\begin{equation}
\label{eq:recursion-hside}
\sum_{\lambda\vdash n} \Gamma^{St}_\lambda(\gamma)e_\lambda=(-1)^{n-n'}q^{(n'n d-n)/2}\sum_{\lambda'\vdash n'} (-1)^{n'-\ell(\lambda')}\Gamma^{St}_{\lambda'}(\gamma'')x_n(d\lambda',q)
\end{equation}

\begin{lemma}
\label{lem:coareatoarea}

We have 
$$q^{n'nd/2-n/2}x_n(d\lambda',q)=q^{\frac{(dn'-1)(en'-1)+n'-1}{2}}\prod_{i=1}^{\ell(\lambda_i')}\left(\sum_{\pi\in\Dyck_{e\lambda_i',d\lambda_i'}^<}q^{-\area(\pi)}h_\pi\right)$$

\end{lemma}

\begin{equation}
\label{eq:recursion-eside}
q^{\frac{(dn'-1)(en'-1)+n'-1}{2}}\omega \varphi_{d/e}|_{q\mapsto q^{-1}} \omega^{-1} (\mathbf{f}_{\gamma''})=q^{\frac{(dn'-1)(en'-1)+n'-1}{2}}\sum_{\lambda'\vdash n/e} \sigma_{\lambda'}(\gamma'')
H_{d,e,\lambda'}^-
\end{equation} 

\begin{theorem}
\label{thm:comparerecursion}
The right-hand sides of Eqs. \eqref{eq:recursion-hside} and \eqref{eq:recursion-eside} are equal up to a sign. More precisely,
$$q^{n'n d/2-n/2}(-1)^{n-n'}\sum_{\lambda'\vdash n'} (-1)^{n'-\ell(\lambda')}\Gamma^{St}_{\lambda'}(\gamma'')x_n(d\lambda',q)=(-1)^{n-n'}q^{\delta_{n'}} \sum_{\lambda'\vdash n'} \sigma_{\lambda'}(\gamma'')H^-_{e,d,\lambda'}$$
where $\delta_{n'}:=\frac{(dn'-1)(en'-1)+n'-1}{2}$  In particular, the left-hand sides are also equal up to the same sign.

\end{theorem}

\begin{proof}
Dividing out the sign and using Lemma \ref{lem:coareatoarea}, the LHS reads $$q^{\delta_{n'}}\sum_{\lambda'\vdash n'} (-1)^{\frac{n}{e}-\ell(\lambda')}\Gamma^{St}_{\lambda'}(\gamma'')\prod_{i=1}^{\ell(\lambda')}\left(\sum_{\pi\in\Dyck_{e\lambda_i',d\lambda_i'}^<}q^{-\area(\pi)}h_\pi\right)$$
and we can divide out $q^{\delta_{n'}}$ on both sides. 

On the other hand, by definition we have 
$$H^-_{e,d,\lambda_i'}=\sum_{\alpha\vDash \lambda_i'}\sum_{\substack{\pi \in \Dyck_{e\lambda_i',d\lambda_i'}\\ \touch{\pi}=\alpha}} q^{-\area(\pi)}h_\pi$$
where $\touch(\pi)=\alpha$ specifies that $\pi$ touches the diagonal at $\alpha$.

Given two arbitrary compositions, denote by $\alpha+\beta$ their concatenation. By sorting, this gives a partition of $|\alpha|+|\beta|$ of length $\ell(\alpha)+\ell(\beta)$. Given a collection $\vec{\alpha}$ of compositions $$\alpha^{(1)}\vDash \lambda_1',\ldots,\alpha^{\ell(\lambda')}\vDash \lambda'_{\ell(\lambda')}$$ refining the parts of $\lambda'$,  we write $\vec{\alpha}\leftrightarrow \lambda'$.  

We now further expand the LHS and collect terms as follows. Note that from the equation $\sum_{k=1}^n(-1)^k h_{n-k}e_k=0$ it follows that 
$$e_n=\sum_{\alpha \vDash n}(-1)^{\ell(\alpha)} h_\alpha$$ where we sum over all compositions of $n$.
Writing $\lambda'=(\lambda_1',\ldots,\lambda_{\ell(\lambda')}')$ we then have 
\begin{equation}
\label{eq:altsum}
e_{\lambda'}=\prod_{i=1}^{\ell(\lambda')}(\sum_{\alpha \vDash \lambda_i'} (-1)^{\ell(\alpha)}h_\alpha)
\end{equation}
By definition of the Dyck germs of $\mathbf{f}_{\gamma''}$, $$\sum_{\lambda'\vdash n/e} \Gamma^{St}_{\lambda'}(\gamma'')e_{\lambda'}=\sum_{\lambda'\vdash n/e} \sigma_{\lambda'}(\gamma'')h_{\lambda'}$$
Fix $\mu'\vdash n/e$. Plugging Eq. \eqref{eq:altsum} in the expression
$$\sum_{\lambda'\vdash n/e} \Gamma^{St}_{\lambda'}(\gamma'')e_{\lambda'}$$
and collecting all collections of compositions whose sum has associated partition $\mu'$ we see that
$$\sigma_{\mu'}(\gamma'')=\sum_{\substack{\vec{\alpha}\leftrightarrow \lambda'\vdash n/e\\\text{sort}(\alpha^{(1)}+\cdots+\alpha^{(k)})=\mu'}}
(-1)^{\frac{n}{e}-\ell(\lambda')}\Gamma^{St}_{\lambda'}(\gamma'')$$ where the sum runs over all $\lambda'\vdash n/e$ and all collections of compositions $\vec{\alpha}$ refining the parts of $\lambda'$, such that the partition given by adding the compositions and sorting is exactly $\mu'$.
It now remains to replace $H^-_{d,e,\lambda'}$ by a similar expansion.
By definition we have 
$$
H^-_{d,e,\lambda'}=\prod_{i=1}^{\ell(\lambda')}
(\sum_{\alpha\vDash \lambda_i'}\sum_{\substack{\pi \in \Dyck_{e\lambda_i',d\lambda_i'}\\ \touch{\pi}=\alpha}} q^{-\area(\pi)}h_\pi)
$$
Picking one of the summands over $\alpha$, we notice that  $$\sum_{\substack{\pi \in \Dyck_{e\lambda_i',d\lambda_i'}\\ \touch{\pi}=\alpha}} q^{-\area(\pi)}h_\pi=
\prod_{i=1}^{\ell(\alpha)}H^{-,<}_{d,e,\alpha_i}$$ 
where we use the fact that the area statistic is additive on concatenation of Dyck paths (but note the coarea is not). Here $H^{-,<}_{d,e,\alpha_i}$ is defined as above.
Again collecting all $\vec{\alpha}\leftrightarrow \mu'$ we see that 
the terms contributing to $H_{d,e,\mu'}^{-}$ on the right are of the form
$$(-1)^{\frac{n}{e}-\ell(\lambda')}\Gamma^{St}_{\lambda'}(\gamma'') \prod_{k=1}^{\ell(\lambda')}\prod_{i=1}^{\ell(\alpha^{(k)})}H^{-,<}_{d,e,\alpha^{(k)}_i}$$
where $\vec{\alpha}$ refines parts of $\lambda'$ and sums and sorts to $\mu'$. Summing over all such collections
we get the desired result.
\end{proof}

\end{document}
